\section{Erd\H{o}s Problem 689 --- Round 2}

\subsection*{1) ROUND-2 OBJECTIVE}

We pursue path \textbf{(C)}: a rigorous obstruction to a major construction family.
Round 1 identified the main unresolved issue: whether a ``first-stage'' choice
\(a_p\equiv 0\pmod p\) for many small primes can be completed by the remaining primes.
In this round I prove that a substantial and natural subclass of this strategy is impossible.
In particular, for the original \(r=2\) problem I rule out every construction with
\[
a_p\equiv 0 \pmod p \qquad\text{for all primes }p\le y
\]
whenever \(y\ge \max(5,n/6)\).
More generally, for the \(r\)-fold variant I rule out all such constructions when \(y\ge n/r\).
This is a genuine advance beyond Round 1: it proves a theorem that excludes an entire class
of candidate constructions that Round 1 had only discussed heuristically.

\subsection*{2) Round-1 FOUNDATION USED}

I use the following vetted results from Round 1.
\begin{enumerate}
\item \textbf{CRT shift equivalence (Round-1 Lemma 4.1).}
Choosing residues \((a_p)_{p\le n}\) is equivalent to choosing a shift \(t\pmod{P(n)}\),
where \(P(n)=\prod_{p\le n}p\), such that
\[
c(m)=\#\{p\le n:\ p\mid (t+m)\}.
\]
Hence the original \(r=2\) problem is equivalent to asking whether there exists a length-\(n\)
block \(t+1,\dots,t+n\) in which every integer has at least two distinct prime divisors \(\le n\).

\item \textbf{Incidence bound (Round-1 Lemma 4.2).}
If every \(m\in[1,n]\) satisfies at least \(r\) congruences, then
\[
\sum_{p\le n}\left\lceil \frac{n}{p}\right\rceil \ge rn.
\]

\item \textbf{All-zero choice lemma (Round-1 Lemma 4.4).}
If \(a_p\equiv 0\pmod p\) for all primes \(p\le n\), then for \(2\le m\le n\) one has
\[
c(m)=\omega(m),
\]
where \(\omega(m)\) is the number of distinct prime divisors of \(m\).
Thus all non-prime-powers in \([2,n]\) are already covered at least twice.

\item \textbf{Computation (Round 1).}
There is no solution for \(n\le 20\) in the original \(r=2\) problem.
\end{enumerate}

No Round-1 proof is repeated here.

\subsection*{3) NEW INSIGHT / TOOL (ROUND-2)}

The new idea is to quantify the \emph{residual demand} created by a partial assignment of residue classes.
A partial assignment already forces each \(m\in[1,n]\) to require a certain number of additional hits.
The remaining primes have only a fixed total incidence budget.
Comparing these two quantities yields a sharp obstruction.

\subsection*{4) ATTACK PLAN (ROUND-2)}

The exact gap left by Round 1 was this: even if one starts with the natural choice
\(a_p\equiv 0\pmod p\) for many small primes, can the remaining primes finish the job?

I will prove two new results.
\begin{enumerate}
\item A \emph{partial-stage capacity lemma} that refines Round-1 Lemma 4.2 and measures the exact remaining demand after any partial choice of residues.
\item Two applications:
  \begin{itemize}
  \item a general \(r\)-fold obstruction showing that the strategy \(a_p\equiv 0\pmod p\) for all \(p\le y\) already fails whenever \(y\ge n/r\);
  \item a sharpened theorem for the original \(r=2\) problem showing that this all-zero strategy already fails whenever \(y\ge n/6\) (more precisely, whenever \(y\ge \max(5,n/6)\)).
  \end{itemize}
\end{enumerate}

The second theorem strictly improves the first when \(r=2\), and it rigorously rules out a major construction family that remained viable after Round 1.

\subsection*{5) WORK (ROUND-2)}

\paragraph{Lemma 5.1 (partial-stage capacity lemma).}
Fix integers \(n\ge 1\) and \(r\ge 1\).
Let \(S\subseteq \mathcal P(n)\) be any set of primes, and suppose residue classes \(a_p\pmod p\) have already been chosen for all \(p\in S\).
For \(m\in\{1,\dots,n\}\), define the current coverage
\[
h_S(m):=\#\{p\in S:\ m\equiv a_p\pmod p\}
\]
and the residual demand
\[
d_S(m):=\max\{0,\,r-h_S(m)\}.
\]
If the partial assignment on \(S\) can be extended to a full choice of residues \((a_p)_{p\le n}\)
such that every \(m\in[1,n]\) satisfies at least \(r\) congruences, then necessarily
\[
\sum_{p\in\mathcal P(n)\setminus S}\left\lceil \frac{n}{p}\right\rceil
\ge
\sum_{m=1}^n d_S(m).
\]

\emph{Proof.}
Assume the partial assignment extends to a full one with final coverage
\[
c(m):=\#\{p\le n:\ m\equiv a_p\pmod p\}\ge r
\qquad (1\le m\le n).
\]
For each remaining prime \(p\in \mathcal P(n)\setminus S\), let
\[
C_p:=\#\{m\in[1,n]:\ m\equiv a_p\pmod p\}.
\]
Then \(C_p\le \lceil n/p\rceil\) for every such \(p\).
Also, for each \(m\in[1,n]\), the number of additional congruences contributed by the remaining primes is
\[
c(m)-h_S(m)\ge r-h_S(m).
\]
If \(h_S(m)\ge r\), then \(d_S(m)=0\), so certainly \(c(m)-h_S(m)\ge d_S(m)\).
If \(h_S(m)<r\), then \(d_S(m)=r-h_S(m)\), and again \(c(m)-h_S(m)\ge d_S(m)\).
Thus in all cases
\[
c(m)-h_S(m)\ge d_S(m).
\]
Summing over \(m\in[1,n]\) gives
\[
\sum_{m=1}^n \bigl(c(m)-h_S(m)\bigr)
\ge
\sum_{m=1}^n d_S(m).
\]
The left-hand side counts exactly the incidences coming from the remaining primes, so
\[
\sum_{m=1}^n \bigl(c(m)-h_S(m)\bigr)=\sum_{p\in\mathcal P(n)\setminus S} C_p.
\]
Hence
\[
\sum_{p\in\mathcal P(n)\setminus S} C_p
\ge
\sum_{m=1}^n d_S(m).
\]
Using \(C_p\le \lceil n/p\rceil\) yields
\[
\sum_{p\in\mathcal P(n)\setminus S}\left\lceil \frac{n}{p}\right\rceil
\ge
\sum_{m=1}^n d_S(m),
\]
as claimed.
\hfill$\square$

\paragraph{Corollary 5.2 (general \(r\)-fold obstruction at threshold \(n/r\)).}
Fix \(r\ge 2\). Let \(y\) be a real number with \(y\ge n/r\).
Suppose one chooses
\[
a_p\equiv 0\pmod p \qquad\text{for every prime }p\le y.
\]
Then there is no extension of these residue classes to all primes \(p\le n\) such that every \(m\in[1,n]\) satisfies at least \(r\) congruences.

\emph{Proof.}
Take \(S:=\{p\in\mathcal P(n): p\le y\}\).
For \(m=1\) one has \(h_S(1)=0\), hence \(d_S(1)=r\).
Also, if \(q\) is any prime with \(y<q\le n\), then \(q\) is not divisible by any prime \(\le y\), so again \(h_S(q)=0\), hence \(d_S(q)=r\).
Therefore
\[
\sum_{m=1}^n d_S(m)
\ge
r\bigl(1+\pi(n)-\pi(y)\bigr).
\]
On the other hand, if \(p>y\), then \(p>n/r\), so
\[
\left\lceil \frac{n}{p}\right\rceil \le r.
\]
Hence the total capacity of the remaining primes is at most
\[
\sum_{p>y}\left\lceil \frac{n}{p}\right\rceil
\le
r\bigl(\pi(n)-\pi(y)\bigr).
\]
This is strictly smaller than \(r(1+\pi(n)-\pi(y))\).
By Lemma 5.1 no extension is possible.
\hfill$\square$

\paragraph{Theorem 5.3 (sharpened \(r=2\) obstruction at threshold \(n/6\)).}
Let \(y\) be a real number with
\[
y\ge \max\left(5,\frac{n}{6}\right).
\]
Suppose one chooses
\[
a_p\equiv 0\pmod p \qquad\text{for every prime }p\le y.
\]
Then there is no extension of these residue classes to all primes \(p\le n\) such that every \(m\in[1,n]\) satisfies at least two congruences.

\emph{Proof.}
Again let
\[
S:=\{p\in\mathcal P(n): p\le y\},
\]
and let \(h_S,d_S\) be as in Lemma 5.1 with \(r=2\).

For each prime \(q\) with \(y<q\le n\), define
\[
U_q:=\{q,2q,3q,4q,5q\}\cap [1,n].
\]
Because \(q>y\ge n/6\), one has \(6q>n\), so
\[
U_q=\{q,2q,\dots,u_q q\},
\qquad u_q:=\left\lfloor\frac{n}{q}\right\rfloor\in\{1,2,3,4,5\}.
\]

I first claim that the sets \(U_q\) are pairwise disjoint as \(q\) ranges over primes \(>y\).
Indeed, suppose
\[
aq=bq'
\]
for some distinct primes \(q,q'>y\) and some integers \(1\le a,b\le 5\).
Without loss of generality, assume \(q'>q\).
Then \(q'\mid aq\).
Since \(q'>q\) and \(q'\) is prime, \(q'\nmid q\).
Also \(q>y\ge 5\), so \(q\ge 7\), hence \(q'>q\ge 7\).
Because \(1\le a\le 5\), we have \(q'>a\), so \(q'\nmid a\).
This is impossible.
Hence \(q=q'\), and then necessarily \(a=b\).
So the \(U_q\) are pairwise disjoint.

Next I compute the residual demand on each \(U_q\).
For the element \(q\in U_q\), no prime \(\le y\) divides \(q\), so
\[
h_S(q)=0,
\qquad d_S(q)=2.
\]
Now let \(aq\in U_q\) with \(a\in\{2,3,4,5\}\).
Since \(q>y\), the prime \(q\) does not contribute to \(h_S(aq)\).
Among the primes \(\le y\), the only prime divisors of \(aq\) are those dividing \(a\).
But each of \(2,3,4,5\) is a prime power, hence has exactly one distinct prime divisor.
Therefore
\[
h_S(aq)=1,
\qquad d_S(aq)=1.
\]
Thus the total residual demand contributed by \(U_q\) is
\[
\sum_{m\in U_q} d_S(m)=2+(u_q-1)=u_q+1.
\]
Since \(u_q=\lfloor n/q\rfloor\), we have
\[
u_q+1\ge \left\lceil\frac{n}{q}\right\rceil.
\]
Therefore
\[
\sum_{m\in U_q} d_S(m)
\ge
\left\lceil\frac{n}{q}\right\rceil.
\]
Because the sets \(U_q\) are pairwise disjoint and \(1\notin U_q\) for any \(q\), we may sum over all primes \(q>y\) and also include \(m=1\), for which \(h_S(1)=0\) and \(d_S(1)=2\). This yields
\[
\sum_{m=1}^n d_S(m)
\ge
2+\sum_{q>y}\left\lceil\frac{n}{q}\right\rceil.
\]
But the total capacity of the remaining primes is at most
\[
\sum_{p>y}\left\lceil\frac{n}{p}\right\rceil.
\]
Hence
\[
\sum_{m=1}^n d_S(m)
>
\sum_{p>y}\left\lceil\frac{n}{p}\right\rceil.
\]
By Lemma 5.1, no extension is possible.
\hfill$\square$

\paragraph{Corollary 5.4.}
Any successful construction for the original \(r=2\) problem with \(n\ge 30\) must use at least one \emph{nonzero} residue among the primes \(p\le n/6\); equivalently, there exists a prime \(p\le n/6\) such that
\[
a_p\not\equiv 0\pmod p.
\]
More generally, in the \(r\)-fold problem any successful construction must use at least one nonzero residue among the primes \(p\le n/r\).

\emph{Proof.}
The \(r=2\) statement follows immediately from Theorem 5.3 once \(n\ge 30\), since then \(n/6\ge 5\).
The general \(r\)-statement follows immediately from Corollary 5.2.
\hfill$\square$

\paragraph{What has been gained relative to Round 1.}
Round 1 only showed that the all-zero construction is a near-solution and identified the difficulty of supplying extra hits to the leftover prime-power-type obstructions.
Round 2 proves a rigorous no-go theorem: one cannot keep \(a_p\equiv 0\pmod p\) for \emph{all} primes up to \(n/6\) (indeed up to \(n/r\) in the \(r\)-fold variant) and hope to complete the construction later.
This rules out an entire class of candidate constructions.

\subsection*{6) ADVERSARIAL VERIFICATION}

\begin{itemize}
\item \textbf{Quantifier check.}
Lemma 5.1 applies to \emph{any} partial assignment on \(S\), not just the all-zero one. Its proof is purely double counting and does not use asymptotics.

\item \textbf{Boundary cases in Corollary 5.2.}
For \(m=1\) and primes \(q>y\), the current coverage is indeed zero when all small residues are \(0\pmod p\).
The strict contradiction comes from the extra target \(m=1\).
No prime number theorem is used.

\item \textbf{Disjointness of the sets \(U_q\).}
This is the only point where the hypothesis \(y\ge 5\) is used.
It ensures all primes \(q>y\) are larger than the multipliers \(1,2,3,4,5\), so no accidental overlaps occur between different \(U_q\).

\item \textbf{Demand computation on \(aq\).}
For \(a\in\{2,3,4,5\}\), the number \(a\) has exactly one distinct prime divisor, so \(aq\) has exactly one prime divisor \(\le y\), namely that divisor of \(a\).
Thus \(d_S(aq)=1\) is correct.

\item \textbf{Why the \(n/6\) threshold appears.}
The proof of Theorem 5.3 uses the prime powers \(2,3,4,5\) as multipliers.
This exactly matches the remaining-prime capacity up to scale \(n/6\).
The same device does not immediately push below \(n/6\), because once one needs multipliers up to \(6\), the set of prime powers is no longer large enough to keep pace with the capacity bound.
So the theorem is sharp for this particular elementary method, though not necessarily sharp for the actual problem.

\item \textbf{Compatibility with Round 1.}
There is no conflict with the Round-1 computation that no full solution exists for \(n\le 20\).
The present theorem is an asymptotic obstruction to a specific family of constructions and is strictly stronger than anything proved in Round 1.
\end{itemize}

\subsection*{7) FINAL}

\noindent\textbf{UNRESOLVED (BUT STRICTLY ADVANCED)}

\smallskip
The original Erd\H{o}s problem remains unresolved here, but Round 2 proves a new theorem that strictly advances beyond Round 1:
\begin{quote}
For the original \(r=2\) problem, any successful construction must depart from the naive choice \(a_p\equiv 0\pmod p\) by the time one reaches primes of size \(n/6\).
More generally, for the \(r\)-fold problem one cannot keep \(a_p\equiv 0\pmod p\) for all primes \(p\le n/r\).
\end{quote}
This rigorously eliminates a substantial class of proof strategies and constructions.

\subsection*{8) COMPLETION ESTIMATE (MANDATORY)}

\noindent\textbf{COMPLETION: 56\%}

\subsection*{9) REFERENCES}

Only references actually used or checked:
\begin{enumerate}
\item P. Erd\H{o}s, \emph{Some problems on number theory}, 1979; the relevant question appears on p.~79.
\item T. F. Bloom, \emph{Erd\H{o}s Problem \#689}, \texttt{https://www.erdosproblems.com/689}, accessed 2026-03-07.
\item B. Green, \emph{100 Open Problems}, Problem 45.
\end{enumerate}
